\section{Introduction}
\label{sec:introduction}

A compound that enters a cell changes what that cell transcribes. Which genes move, by how much,
and in which direction is the most direct readout available of what the compound actually does ---
not what it was designed to bind, but what the cell does in response. Reading that response is
routine. Reading it for every combination that might matter is not. A compound behaves differently
at different doses, and differently again depending on the genetic background of the cell receiving
it, so the space to be explored is the product of three large sets rather than any one of them. Even
the largest screens sample a small corner of it, and every combination left untested is a question
left open.

This is the gap prediction is meant to close. If the transcriptional response to a treatment could
be computed rather than measured, the untested corners of that space would become inspectable, and
the experiments actually run could be chosen because they are informative rather than because they
happened to fit on the plate. The task, stated plainly, is this: given the transcriptome of an
untreated cell and the identity of a drug, predict the transcriptome that cell would have shown had
the drug been applied.

The task is harder than it sounds, for a reason that is easy to state and impossible to engineer
around. The ground truth is a counterfactual. Sequencing destroys the cell it reads, so the treated
and untreated profiles of the same cell can never both be observed, and a prediction can only ever
be checked against a different cell that received the drug. Everything downstream inherits this. It
is why paired training data must be assembled by matching cells rather than by tracking them, why
evaluation is usually framed at the level of populations rather than individuals, and why
distinguishing a real drug effect from ordinary cell-to-cell variation requires care rather than
optimism.

Assembling data at the scale this problem demands became possible only recently. Two decades of
falling sequencing costs, rising throughput and increasingly clever sample multiplexing turned
single-cell profiling from a delicate experiment into something that can be run industrially, and
Tahoe-100M~\cite{tahoe100m} is what the end of that trajectory looks like: roughly a hundred million
single-cell transcriptomes spanning some 1,100 small molecules across 50 cancer cell lines. Every
compound and every cell line appears not as a summary statistic but as a cloud of individual cells.

That resolution changes which questions can be asked. The obvious one is conditional --- what does
this compound do to this cellular context, rather than what does it do on average across contexts
that may have little in common. The less obvious question is whether the first has an answer.
Because each treated condition is a population of cells rather than a single measurement, it becomes
possible to measure how much of an apparent drug effect reproduces across disjoint halves of the
same treated well and how much is only variation between the cells that happened to be sampled. That second measurement turns out to matter as much as the first, and a good deal of what
follows rests on the observation that standard evaluations quietly confuse the two.

The other development behind this thesis has nothing to do with cells. Over the same period,
pretrained transformer language models became unreasonably good and unreasonably general, and the
engineering poured into them --- architectures, training recipes, scaling laws, the entire
surrounding toolchain --- came to exceed what any single scientific field could plausibly assemble
for itself. That created a standing temptation: if a problem can be made to look like text, it
inherits all of that work for free.

Cell2Sentence~\cite{c2s} takes the temptation about as literally as it can be taken. Rank the genes
in a cell from most to least expressed, write their names out in that order, and the cell becomes a
sentence --- ordinary text that an ordinary language model can be fine-tuned on. The encoding is
unsentimental about biology, discarding how much each gene is expressed and keeping only the
ordering, on the argument that rank retains enough of the signal for the rest to be reconstructed.
In exchange it requires no new architecture and no new training machinery, only cells written in a
format that existing models already consume.

C2S-Scale~\cite{c2sscale} develops that idea into a family of models. Larger backbones are
pretrained on over 50 million cells drawn from 825 single-cell datasets in CELLxGENE and the
Human Cell Atlas, and then fine-tuned across a spread of tasks: predicting cell type and tissue,
generating cells conditioned on a label, captioning clusters, answering questions about a dataset in
natural language, and predicting perturbation responses. That last task is the one this thesis
builds on, and it is worth being specific about its scope. Their perturbation work covers cytokine
stimulation in peripheral blood cells, and, for small molecules, bulk L1000 profiles and the
SciPlex3 panel. The single-cell results are scored as distributions --- how closely the generated
population for a condition resembles the observed one --- rather than condition by condition. The
checkpoint used throughout this thesis, \texttt{C2S-Scale-Pythia-1b-pt}, is the pretrained model
from that family: a Pythia-1b language model~\cite{pythia}, pretrained on natural language and then
further pretrained on atlas cells, with no drug-treated cells among that atlas data.

What we do is extend that line of work to chemical perturbation at the scale Tahoe-100M now makes
available, and evaluate it as a prediction problem in the strict sense --- condition by condition,
against controls chosen so that a model which ignores the drug cannot pass. The question is whether
a cell-sentence model fine-tuned on this atlas learns what a particular compound does to a
particular cell.

Three things are worth keeping straight about what the model has actually seen, since the scale of
the atlas invites an easy overstatement. Its cell pretraining is large but drug-free: tens of
millions of cells, none of them treated with anything. The fine-tuning does not consume Tahoe
entire --- it draws a sample of a few hundred thousand treated cells from it, which is ample for the
task but is not a hundred million. And nothing here is unseen by the whole pipeline: the backbone is
the Pythia-1b decoder~\cite{pythia}, a general-purpose language model pretrained on natural text,
and every prompt names the compound outright and carries a \texttt{Mechanism} field --- its
annotated mechanism of action where one exists, the literal string \texttt{unclear} where it does
not. What a held-out drug is held out of is the Tahoe fine-tuning, with its name and whatever
mechanism annotation it carries still supplied at inference --- not the model's prior exposure to
that name.

It does learn something, and the something is the difficulty. The model acquires a good account of
what a treated cell looks like \emph{in general} --- the shared shape that perturbed transcriptomes
have in common, which is substantial and largely the same whichever compound produced it. What it
does not acquire is the part that distinguishes one compound from another. Change the drug named in
the prompt while holding the input cell fixed, and the prediction barely moves; on the calibrated
metric the model is statistically indistinguishable from a baseline that copies the untreated cell
and is told nothing about any drug at all.

A second comparison is worth separating from that one, because it is weaker but more informative. A
per-drug average, estimated from the same compound in \emph{other} cell lines, also outperforms the
model --- but this is not a naive baseline. On the conditions where both are defined ($n = 1192$) it
scores $0.913$ where the model scores $0.549$, so matching it would be a real
achievement rather than a low bar. What it cannot do is condition on the cell
in front of it, which is precisely why the comparison is useful: it separates the drug's average
effect from the part that depends on context. Section~\ref{sec:res-transfer} measures that split with a
transfer coefficient \ci{\Tcoef = 0.553}{0.514}{0.593}, leaving \ci{44.7\%}{40.7\%}{48.6\%} of the
drug-specific residual not shared across contexts. That residual is not a drug-by-cell-line
interaction share --- the estimator mixes cell line, dose, well and target estimator --- and whether
any model could reach it is a separate question, taken up there.

A negative result earns its place only if it can be made specific, and that is the organising
principle of what follows. We establish first that the failure is real rather than an artifact of
how this field measures success --- a distinction less pedantic than it sounds, since the standard
metric can be saturated by a predictor that knows nothing about drugs. We then locate where in the
pipeline the failure originates, show that the location implies a repair, apply it, and finally
measure how much of the reproducible drug response is not shared across contexts --- across cell
line, dose and well together --- a descriptive quantity rather than a bound on what any method
could reach. The thesis therefore ends not in closure but in a number: a measured fraction of the
drug-specific residual that is not shared across contexts --- a component the per-drug lookup cannot
express by construction, and one no model in this thesis was shown to recover.

Four questions carry that argument, in order:

\begin{enumerate}
  \item \textbf{Is the evaluation trustworthy?} Before asking whether a model is any good, one has
        to ask whether the metric could tell. Can a predictor carrying no drug information score
        well on the measures the field has standardised on, and if so, which measures survive?
  \item \textbf{Does the model use the drug at all?} Holding the input cell fixed and changing only
        the compound it is told about, does the prediction move --- and does it move in the
        direction the real response moved?
  \item \textbf{If not, where does the failure live?} Is the drug missing from the data, missing
        from the model's internal representation, or present in that representation but ignored by
        the process that writes the output? Each answer implies a different repair, and only one of
        them is cheap.
  \item \textbf{How much is recoverable at all?} A per-drug average is a strong baseline. How much
        of the drug-specific signal does it capture, how much of it is not shared across contexts,
        and how much of that unshared part do our models or any baseline we could construct
        actually reach?
\end{enumerate}

The contributions follow the same order. We audit the calibration of the standard evaluation and
show that it can be saturated by a predictor carrying no drug information whatsoever, which bears on
any work built from the same family of metrics. Of the five metrics we test, only the normalised
inverse rank has a positive dynamic range fraction --- $+0.635$ across plates and $+0.545$ within
plate, against $-0.16$ to $-0.69$ for the other four across plates and $-0.10$ to $-0.52$ within
plate. Both are measured with the bootstrap clustered on cell lines rather than on the drug rows
nested in them --- $[+0.604, +0.663]$ across plates over $25$ lines, $[+0.516, +0.569]$ within plate
over $27$, each with Holm $p = 0.002$. We
establish drug-blindness on that one surviving metric, using several independent instruments so that
no single measurement carries the claim alone. We identify the target encoding --- not the data, the architecture, or the training
objective --- as the binding constraint, and derive from that diagnosis a re-encoding of the
prediction target which produces the first measurable drug effect in this project. On training
conditions the model beats its designated comparator (\texttt{scramble\_orth}: the same prompt with
the drug name replaced by a within-plate partner chosen so that its measured response is
near-orthogonal to the target drug's) by $+0.0898$, with an interval excluding zero whether the
resampling is clustered on cell line $[+0.0382,\,+0.1415]$ or on drug $[+0.0291,\,+0.1506]$. We do
\emph{not} claim that this transfers. On held-out drug-and-cell-line pairings the same gap is
$+0.0332$, which excludes zero clustered on cell line $[+0.0069,\,+0.0594]$ but not clustered on
drug $[-0.0031,\,+0.0695]$ --- and drug is the axis across which a claim about transferring a drug
signature has to generalise. That split is in any case mostly recombination rather than novelty:
$30$ of its $34$ drugs and $39$ of its $42$ cell lines also appear in the training split, and $82\%$
of its conditions pair a training drug with a training cell line. On genuinely unseen drugs the gap
is $-0.0315$, with intervals spanning zero under both clusterings ($[-0.0836,\,+0.0206]$ on cell
line, $[-0.1049,\,+0.0419]$ on drug). Nor are the three gaps on a common scale: the within-well
split-half precision reference differs across the splits ($0.956$, $0.824$, $0.836$), so their
ordering cannot be read as a clean decay with novelty. Finally we separate the part of the
reproducible drug response that is shared across contexts from the part that is not, and report the
unshared fraction as $0.447$ $[0.407,\,0.486]$ --- a residual mixing cell line, dose, well and
estimator rather than a clean drug-by-cell-line interaction.

The thesis proceeds as follows. Section~\ref{sec:litreview} places the work against single-cell
foundation models, the perturbation-prediction literature, and an unsettled dispute within it about
whether deep models beat trivial baselines at all --- a dispute this thesis adjudicates in one
corner, using the instruments of the side it ultimately argues against. Section~\ref{sec:methods}
is the investigation itself, told in the order it happened: because nearly every instrument in it
was built in response to the measurement that came before, the data, the methods and the findings
are presented as one continuous account, with each definition given where it is first needed.
Section~\ref{sec:limitations} sets out the scope honestly, separating what is closed by measurement
from what is merely untested, and Section~\ref{sec:conclusions} concludes.
